\section{Methods}
\label{sec:methods}

\subsection{Robots and simulation}
\label{sec:robots}

Robots are built from three module types: a core (half-extent \SI{0.075}{m}, mass \SI{1.363}{kg}), passive bricks (half-extent \SI{0.0375}{m}, \SI{0.055}{kg}) and actuated hinges (stator \SI{0.065}{kg}, rotor \SI{0.040}{kg}). Each hinge is a position servo with joint limits of $\pm 90^\circ$, gain $k_p = 1.36$, damping $k_v = 0.356$ and armature $0.081$, values fitted to measurements of the physical hinge module.
\todo{Add the source of the hinge fit, and a figure with example bodies of each genome.}

We simulate with MuJoCo on a flat floor, with physics at \SI{500}{Hz} and control at \SI{20}{Hz}. A body spawns with its core at \SI{0.1}{m} and settles for \SI{3}{s} with zero control before each episode.

\subsection{Body genomes}
\label{sec:genomes}

All bodies are encoded by a tree in which each node is a module, with the face of its parent it attaches to and a rotation about that face (multiples of $45^\circ$). New bodies have at most 25 modules plus the core, a tree depth of at most 25, and at least four hinges. We compare three genomes.

The tree genome is unconstrained. The symmetric genome mirrors every subtree about the vertical plane through the core that contains the forward axis, and forces midline modules to rotation $0^\circ$. The tied genome duplicates without mirroring: a placed module receives, with probability $p = 0.47$, an identical twin at a random free site. This makes 64\% of modules part of a tied pair, the same share as the mirrored modules in random symmetric bodies.

Each genome has its own random initialisation, subtree crossover (probability 0.6) and mutation operators, which preserve its constraint. The module limit does not hold strictly under these operators: evolved tied bodies exceed it often (42\% of the final population under \bb, with up to 64 modules across controllers) and tree bodies rarely (about 2\%). Fixed-body experiments are not affected.

\subsection{Task and fitness}
\label{sec:task}

The task is forward locomotion. An episode lasts \SI{15}{s} after the settle. Fitness is the displacement of the core along the forward axis in metres, without clipping or penalties; moving backwards gives a negative score. The simulation is deterministic on one machine, so each controller is evaluated once.

\subsection{Controllers}
\label{sec:brains}

We use two controllers, plus two variants of the first. The \bb\ controller is open-loop: each hinge follows a square wave between $-90^\circ$ and $+90^\circ$, all hinges share one frequency (\SIrange{0.2}{2}{Hz}), and each hinge has its own duty cycle and phase, which gives $4 + 5n$ parameters for $n$ hinges. \Bbfree\ gives every hinge its own frequency ($9n$ parameters). \Bbdrift\ has as many parameters but keeps the frequencies within a small offset of the shared one ($4 + 9n$), which separates losing the shared clock from enlarging the search space.

\dann\ is a closed-loop decentralised network. One multilayer perceptron (42 inputs, 16 hidden units, one output; 705 parameters) is shared by every hinge. Each hinge sees its own angle, velocity and module type, the same for its nearest hinge in each of six directions, a sine clock, the core orientation and its own rest position.

\subsection{Learning and evolution}
\label{sec:learning}

Each body's controller is learned by CMA-ES ($\sigma_0 = 0.5$, population 20) from a fixed starting mean, and a body's score is the best displacement found.

We run two experiments (\cref{tab:experiments}). In the evolution experiment, a $(\mu + \lambda)$ evolutionary algorithm with $\mu = \lambda = 20$, tournament selection (size 4) and plus-selection runs for 30 generations, with 200 controller evaluations per body. Bodies are trained once and parents are not retrained. A run's score is the best body in its final population.

In the fixed-body experiment, controllers are trained without evolution: 1{,}600 evaluations for the \bb\ family and 5{,}000 for \dann, with three training seeds per body. There are 80 bodies per set, and all sets have the same hinge counts. Tree and symmetric bodies are drawn from each genome's initial distribution and paired on hinge count. Tied bodies are drawn the same way. Shape-matched tree bodies are selected from 24{,}000 random tree bodies, one per symmetric body, with the closest resting shape on six standardised measures (leg hinges, yaw-axis leg hinges, front--back contact spread, leg reach, height, lateral-axis hinges), assigned by the Hungarian algorithm within each hinge count. We chose these measures because they predicted score in earlier pilot runs. Leg-matched tree bodies are selected in the same way to match the tied bodies on leg hinges.

\begin{table}[t]
  \centering
  \small
  \caption{Experiments. Evaluations are controller evaluations per body.}
  \label{tab:experiments}
  \resizebox{\linewidth}{!}{\begin{tabular}{llll}
    \toprule
    experiment & bodies & controllers & evaluations \\
    \midrule
    evolution, 20 seeds & tree, symmetric, tied & \bb, \bbdrift, \bbfree & 200 \\
    fixed bodies, 3 seeds per body & 80 tree, 80 symmetric & all four & 1{,}600 (\dann\ 5{,}000) \\
    & 80 tied, 80 shape-matched, 80 leg-matched & \bb, \dann & 1{,}600 (\dann\ 5{,}000) \\
    \bottomrule
  \end{tabular}}
\end{table}

\subsection{Body and gait measures}
\label{sec:measures}

We measure the resting shape after the settle. Height is the vertical extent of the module positions, and contact spread is the standard deviation of floor contacts along the forward axis. Each hinge is classed by the dominant world direction of its rotation axis: lateral (the limb swings fore and aft), forward, or vertical (yaw, the limb sweeps horizontally). A leg hinge is a hinge whose distal part touches the floor, and leg reach is the mean distance from a leg hinge to its distal floor contacts.

To measure gaits, we replay trained controllers and pair left and right modules by matching rest positions across the body's sagittal plane; this procedure is the same for every genome. For each pair, the trajectories are cut into one-period windows and the phase lag is computed per window. Phase quality is one minus the circular variance of those lags, so 1 means the same lag in every cycle. A pair is locked when its phase quality is at least 0.9. A gait's locked share is the share of its stroke carried by locked pairs, and a gait is locked when this share is at least 0.5. Because pairs come from near-mirrored modules, tree and tied bodies have fewer pairs by construction.
\todo{Watch about 20 replays per group to confirm that the detector holds on these modules.}

\subsection{Phase-shift replay}
\label{sec:phase-shift}

We replay each trained controller from the fixed-body experiment with the timing of its left-side hinges perturbed. A phase shift adds a constant offset of 0.05, 0.1, 0.25 or 0.5 cycles, which changes the lag to the other hinges but keeps it constant. Detuning runs the hinges at $1 + D$ times the shared frequency, with $D = 0.02$ or $0.15$, so the lag drifts and the lock breaks. For \dann\ the perturbation acts on the hinges' clock input, and the closed loop can partly compensate.

As a control, the same perturbation is applied to the same number of randomly chosen lateral hinges (three draws). Shifting any $k$ of $n$ hinges changes the relative timing of $k(n-k)$ hinge pairs, whichever hinges are chosen, so the two conditions disrupt equally many relations. They differ in that the side perturbation disrupts exactly the left--right ones.

\subsection{Statistics}
\label{sec:stats}

Evolution runs are paired by seed; we report the mean paired difference, a bootstrap 95\% CI over seeds and a Wilcoxon signed-rank test. On fixed bodies, seeds are averaged per body, so the unit is the body. We report the difference of means, a bootstrap 95\% CI over bodies, a permutation test on the mean and a Mann--Whitney $U$ test. We report both tests because scores are skewed, with a few bodies far ahead of the rest, so tests of means and ranks can disagree. Phase-shift differences are averaged per body and tested with a Wilcoxon test. Tests are not corrected for multiple comparisons, and the analyses of body measures are exploratory.
